\section{15 октября 1914 г.}

\lp{20260604_193210861}

Дорогой мой, милый!

Был у Коли в гостях Сабуров и сказал, чтобы я вас непременно предупредила, что
Швейцария может быть тоже втянута, и тогда положение русских очень трудное,
почти безвыходное, т.к. могут закрыться и те пути, которые сейчас есть.

\section{16 октября 1914 г.}

\lp{20260604_193210861}

Милый! Сегодня у нас были Киселевы и папаша. Часть вечера (в начале) пробыл и
Саша Гедике. Он с нескрываемым восхищением смотрит на эту великолепную
Валькирию. В 10 часов ушел и папаша, и мы просидели с Киселевыми до 11.
Говорилось, конечно, опять о событиях. Причем Коля опять волновался. Он все
время страшно волнуется и меня это очень беспокоит. Он никак не может принять
войны. На него больно смотреть. Он так много говорит, что сделался даже
красноречив. Работа его не идет. Что то будет, что то будет\textellipsis

Сейчас иду спать, а завтра продолжу. Посылаю вам письмо для Мариэтты: не знаю,
где

\lp{20260604_193219210}

\noindent
она сейчас. Уехала ли в Италию как собиралась или может быть даже и вернулась
оттуда.

Киселев оставил мне список самых необходимых Мусагетских расходов, и теперь я
могу поговорить с Марго, т.к. дело выяснилось, и я даже могу ей показать точный
список расходов.

Покойной ночи, мой дорогой!

\section{17 октября 1914 г.}

\lp{20260604_193219210}

Сейчас пришло еще одно ваше заказное шестое \rem{письмо} с великолепной головой
очаровательного старика. Но вы все-таки этого не делайте и письма пишите более
отвлеченные. Я уж догадаюсь что и о чем вы хотите сказать. Вы ведь и не
представляете себе, сидя в нейтральной стране, до чего здесь такое письмо может
произвести неприятное впечатление и вызвать подозрение. Имейте ввиду, что
действует не только военная цензура, которая безусловно

\lp{20260604_193219210}

\noindent
ведет себя благородно, но и тайная полиция, которая не очень то станет
разбирать и решит, что высказанное положительное по отношении к одной стороне
есть уже несомненно отрицательное по отношении к другой.

Письма ваши все дошли, это ясно видно по связи и по числам. Швейцарка тоже
привезла мне ваше письмо и открытку от Мариэтты.

Начинаю беспокоиться. Когда же вы будете дома? А то, если весь мир будет
захвачен войной? И сколько времени понадобится, чтобы погасить такое пламя?

Сейчас была опять Heimbach. Просидели довольно долго. А теперь мне уже и опять
надо торопиться. Мы выйдем сейчас с Колей вместе, зайдем в Гнездниковский
\rem{переулок}, там поужинаем; оттуда Коля пойдет к Рахманинову, а я домой, и
тогда еще попишу вам, мой дорогой, милый!

\lp{20260604_193210861}

Не знаю написал ли вам уже Ильин, спрошу его об этом.

Это письмо я все-таки недописанное возьму с собой и папаша его рано утром
отправит.

Третьего дня я опустила письмо с маленькой Колиной припиской. Неужели оно не
дошло?

Ваша Анюта.
